\begin{recipeDP}
    [
        preparationtime = {\SI{15}{\minute}},
        bakingtime = {\SI{60}{\minute} bis \SI{65}{\minute}},
        bakingtemperature = {\protect\bakingtemperature{fanoven=\SI{180}{\celsius}}},
        portion = {1 Kastenform},
        source = {lillakocht.de}
    ]
    {Karottenkuchen}

    \graph
        {
            small=Gebaeck/Kuchen/Karottenkuchen/small.jpg,
            big=Gebaeck/Kuchen/Karottenkuchen/big.jpg
        }

    % \introduction{
    %     Ein leichter, saftiger veganer Karottenkuchen mit warmen Gewürzen, perfekt für Kaffee oder als süßer Snack.
    % }

    \ingredients{
        \multicolumn{2}{l}{\textbf{Teig}} \\
        \SI{200}{\g} & Mehl \\
        \SI{80}{\g} & Zucker \\
        \SI{80}{\g} & brauner Zucker \\
        \SI{1}{\pck} & Backpulver \\
        \SI{2}{\TL} & Zimt \\
        \SI{1}{\TL} & Muskatnuss \\
        \SI{1,5}{\TL} & Spekulatiusgewürz \\
        \SI{1}{\TL} & Ingwer, gemahlen \\
        \SI{1}{\TL} & Salz \\
        \SI{300}{\g} & grob geriebene \addtoidx{Karotten} \\
        \SI{150}{\ml} & Pflanzenmilch \\
        \SI{100}{\ml} & Öl \\
        \SI{100}{\g} & Walnüsse \\
        \SI{1}{\EL} & Apfelessig \\
        \SI{1}{\TL} & Vanilleextrakt \\
        \\
        \multicolumn{2}{l}{\textbf{Frosting}} \\
        \SI{125}{\g} & Butter \\
        \SI{250}{\g} & Frischkäse \\
        \SI{200}{\g} & Puderzucker \\
        \SI{1}{\TL} & Vanilleextrakt \\
        \SI{50}{\g} & Walnüsse
    }

    \preparation{
        \step Den Ofen auf \SI{180}{\celsius} vorheizen und eine Kastenform ( \SI{20}{\cm}) mit Backpapier auslegen.

        \step In einer großen Schüssel Mehl, Zucker, braunen Zucker, Backpulver, Zimt, Muskatnuss, Ingwer und Salz vermengen.
        In einer zweiten Schüssel die geriebenen Karotten mit Pflanzenmilch, Öl, Nüssen, Apfelessig und Vanilleextrakt verrühren.

        \step Die feuchten Zutaten in die trockenen geben und nur so lange rühren, bis alles gerade vermengt ist.
        Der Teig sollte leicht klumpig sein.
        Den Teig gleichmäßig in die Form füllen und leicht aufklopfen, damit keine Luftblasen entstehen.

        \step Den Kuchen etwa \SI{60}{\minute} bis \SI{65}{\minute} backen, bis ein Zahnstocher sauber herauskommt.
        Danach \SI{10}{\minute} in der Form abkühlen lassen und anschließend aus der Form stürzen.

        \step Für das Frosting die vegane Butter weich rühren, dann Frischkäse, Puderzucker und Vanilleextrakt hinzufügen und cremig mixen.
        Den vollständig abgekühlten Karottenkuchen mit dem Frosting bestreichen und nach Belieben mit gehackten Nüssen dekorieren.
    }

    % \hint{
    %     Rosinen, Kokosflocken oder Kürbiskerne passen gut in den Teig.
    % }

\end{recipeDP}
